% ╔══════════════════════════════════════════════════════════════╗
% ║              功能筆記 README — Template_Slides               ║
% ║  說明本模板的使用方式、資料夾結構與自訂環境                  ║
% ║  本檔不會被編譯，僅供參考                                    ║
% ╚══════════════════════════════════════════════════════════════╝


% ────────────────────────────────────────────────────────────
%  一、資料夾結構
% ────────────────────────────────────────────────────────────
%
%  Template_Slides/
%  ├── note-template/                ← 模板核心（請勿隨意修改）
%  │   ├── main.tex                  ← 編譯入口，切換文件只改這裡
%  │   ├── package_and_settings.tex  ← 套件、色彩主題、自訂環境
%  │   ├── 功能筆記readme.tex        ← 本說明文件
%  │   └── fonts/                    ← 字型檔
%  │
%  └── Content/                      ← 所有文件內容
%      ├── course/                   ← 課堂講義
%      │   ├── {科目1}/
%      │   │   ├── note/note_1/      ← 上課投影片
%      │   │   ├── hw/hw_1/          ← 作業講解
%      │   │   ├── quiz/quiz_1/      ← 小考
%      │   │   └── exam/midterm|final/
%      │   └── {科目2}/         （結構同上）
%      │
%      ├── paper/                    ← 論文報告（讀書會、課堂）
%      │   ├── {作者_年份1}/
%      │   └── {作者_年份2}/
%      │
%      ├── thesis/                   ← 論文口試簡報
%      │   └── {論文口試}/
%      │
%      └── meetings/                 ← 會議報告
%          └── {對象}/       ← 按報告對象建子資料夾
%              └── {YYYYMMDD}/         ← 日期格式：YYYYMMDD
%
%  ► 新增會議：在 meetings/{對象}/ 下建立新的日期資料夾，複製相同結構。
%  ► 新增論文：在 paper/ 下建立 {作者_年份}/ 資料夾。
%  ► 新增科目：在 course/ 下建立科目資料夾，結構與上方相同。


% ────────────────────────────────────────────────────────────
%  二、每個文件資料夾的內部結構
% ────────────────────────────────────────────────────────────
%
%  {文件資料夾}/
%  ├── 0. 標題與作者.tex  ← 定義 \title / \author / \date
%  ├── part0.tex           ← 標題頁（\frame{\titlepage}）與目錄
%  ├── part0_all.tex       ← 依序 \input 所有 part 檔的索引
%  ├── part1.tex           ← 第一段投影片
%  ├── part2.tex           ← 第二段（視需要新增）
%  └── figures/            ← 本文件用到的圖片


% ────────────────────────────────────────────────────────────
%  三、色彩主題切換
% ────────────────────────────────────────────────────────────
%
%  只需修改 package_and_settings.tex 第一個設定：
%
%    \newcommand{\slidetheme}{green}   % 可改為 blue 或 purple
%
%  三種色調說明：
%    green  ── 深青綠標題列（預設，偏沉穩）
%    blue   ── 深藍標題列（學術感強）
%    purple ── 深紫標題列（個性鮮明）
%
%  改完重新編譯，標題列、頁尾、框線、連結顏色全部跟著換。
%  文字顏色（黑）與警示色（紅）不受影響，確保可讀性。


% ────────────────────────────────────────────────────────────
%  四、如何切換文件（main.tex 的操作）
% ────────────────────────────────────────────────────────────
%
%  只需修改 main.tex 中的兩行 \input：
%
%    步驟一：\input{Content/{類型}/{...}/0. 標題與作者}
%    步驟二：\input{Content/{類型}/{...}/part0_all}
%
%  範例：
%    切換到會議報告（{對象}，{日期}）：
%      \input{Content/meetings/{對象}/{YYYYMMDD}/0. 標題與作者}
%      \input{Content/meetings/{對象}/{YYYYMMDD}/part0_all}
%
%    切換到 {作者 et al.}（{年份}）論文報告：
%      \input{Content/paper/{作者_年份1}/0. 標題與作者}
%      \input{Content/paper/{作者_年份1}/part0_all}
%
%    切換到碩士論文口試：
%      \input{Content/thesis/{論文口試}/0. 標題與作者}
%      \input{Content/thesis/{論文口試}/part0_all}


% ────────────────────────────────────────────────────────────
%  五、圖片管理
% ────────────────────────────────────────────────────────────
%
%  圖片放在各文件的 figures/ 子資料夾：
%    Content/meetings/{對象}/{YYYYMMDD}/figures/my_figure.png
%
%  引用時使用從根目錄出發的完整相對路徑：
%
%    \begin{frame}{含圖投影片}
%      \centering
%      \includegraphics[width=0.7\textwidth]
%        {Content/meetings/{對象}/{YYYYMMDD}/figures/my_figure}
%    \end{frame}
%
%  注意：不要只寫檔名，Overleaf 從 main.tex 所在的資料夾出發找路徑。


% ────────────────────────────────────────────────────────────
%  六、自訂環境使用說明
% ────────────────────────────────────────────────────────────

% ── problembox（題目框）──────────────────────────────────────
%
%  \begin{problembox}{題號}{題目名稱}
%    題目內容...
%  \end{problembox}
%
%  ► 若題目很長，需要跨多張投影片：
%    \begin{frame}[allowframebreaks]{標題}
%      \begin{problembox}{1.1}{題目名稱}
%        非常長的題目內容...
%      \end{problembox}
%    \end{frame}

% ── solutionbox（解答框）────────────────────────────────────
%
%  \begin{solutionbox}
%    解題過程...
%  \end{solutionbox}
%
%  ► problembox + solutionbox 的標準搭配：
%    \begin{frame}[allowframebreaks]{題目標題}
%      \begin{problembox}{1.1}{OLS 估計}
%        給定模型 $y = X\beta + \epsilon$，推導 OLS 估計量...
%      \end{problembox}
%      \begin{solutionbox}
%        最小化 $\|y - X\beta\|^2$，對 $\beta$ 偏微分得...
%        \[\hat\beta = (X^\top X)^{-1}X^\top y\]
%      \end{solutionbox}
%    \end{frame}

% ── highlightbox（強調框，可自訂標題）──────────────────────
%
%  \begin{highlightbox}{框標題}
%    重要定理或結論...
%  \end{highlightbox}

% ── Beamer 內建 block 環境（顏色自動跟隨主題）──────────────
%
%  \begin{block}{一般說明框}
%    內容...
%  \end{block}
%
%  \begin{exampleblock}{範例框}
%    範例內容...
%  \end{exampleblock}
%
%  \begin{alertblock}{警示框}
%    重要警告...
%  \end{alertblock}

% ── \alert{} 行內標記（變成警示色）─────────────────────────
%
%  \alert{這段文字會變成紅色}

% ── 其他排版工具 ──────────────────────────────────────────
%
%  多欄排版：
%    \begin{columns}
%      \begin{column}{0.48\textwidth}
%        左欄內容...
%      \end{column}
%      \begin{column}{0.48\textwidth}
%        右欄內容...
%      \end{column}
%    \end{columns}
%
%  表格：推薦使用 booktabs 樣式
%    \begin{tabular}{lcc}
%      \toprule
%      欄位 A & 欄位 B & 欄位 C \\
%      \midrule
%      資料 & 資料 & 資料 \\
%      \bottomrule
%    \end{tabular}
%
%  字型指令：
%    \kai{楷體}   \ming{明體}   \hei{黑體}

